\documentclass[11pt]{article}
\usepackage[margin=1in]{geometry}
\usepackage{amsmath,amssymb,amsthm,mathtools}
\usepackage{microtype}
\input{glyphtounicode}
\pdfgentounicode=1
\pdfglyphtounicode{mu}{03BC}
\pdfglyphtounicode{bardbl}{2016}
\pdfglyphtounicode{parenleftbig}{0028}
\pdfglyphtounicode{parenrightbig}{0029}
\pdfglyphtounicode{parenleftbigg}{0028}
\pdfglyphtounicode{parenrightbigg}{0029}
\pdfglyphtounicode{parenlefttp}{239B}
\pdfglyphtounicode{parenleftbt}{239D}
\pdfglyphtounicode{parenrighttp}{239E}
\pdfglyphtounicode{parenrightbt}{23A0}
\pdfglyphtounicode{circleplusdisplay}{2A01}
\pdfglyphtounicode{summationdisplay}{2211}
\pdfglyphtounicode{productdisplay}{220F}
\usepackage[colorlinks=true,linkcolor=blue,citecolor=blue,urlcolor=blue]{hyperref}
\usepackage{accsupp}
\NewCommandCopy{\OriginalHookrightarrow}{\hookrightarrow}
\NewCommandCopy{\OriginalLongmapsto}{\longmapsto}
\NewCommandCopy{\OriginalUpbracefill}{\upbracefill}
\renewcommand{\hookrightarrow}{\mathrel{%
  \BeginAccSupp{method=hex,unicode,ActualText=21AA}%
  \OriginalHookrightarrow\EndAccSupp{}}}
\renewcommand{\longmapsto}{\mathrel{%
  \BeginAccSupp{method=hex,unicode,ActualText=27FC}%
  \OriginalLongmapsto\EndAccSupp{}}}
\renewcommand{\upbracefill}{%
  \BeginAccSupp{method=hex,unicode,ActualText=23DF}%
  \OriginalUpbracefill\EndAccSupp{}}
\hypersetup{pdftitle={Quadratic stabilization of the canonical Foulkes--Howe map},pdfauthor={OpenAI}}
\newtheorem{theorem}{Theorem}
\newtheorem{corollary}[theorem]{Corollary}
\newtheorem{proposition}[theorem]{Proposition}
\newtheorem{lemma}[theorem]{Lemma}
\theoremstyle{remark}
\newtheorem{remark}[theorem]{Remark}
\DeclareMathOperator{\Sym}{Sym}
\DeclareMathOperator{\Av}{Av}
\newcommand{\C}{\mathbb C}
\newcommand{\Sa}{\mathfrak S_a}
\newcommand{\D}[2]{D_{#1,#2}}
\title{Quadratic stabilization of the canonical Foulkes--Howe map}
\author{OpenAI}
\date{September 25, 2026}
\begin{document}
\maketitle
\begin{abstract}
For every finite-dimensional complex vector space $V$, we prove that the
canonical Foulkes--Howe map
$\Sym^b(\Sym^a V)\longrightarrow\Sym^a(\Sym^b V)$ is surjective whenever
$a\ge2$ and $b\ge a(a-1)$. This gives a quadratic stabilization bound
independent of $\dim V$.
\end{abstract}

\section{Introduction}\label{sec:introduction}

Let $V$ be a finite-dimensional complex vector space. Products of vectors
will always mean products in the symmetric algebra $\Sym V$.
For an integer $a\ge1$ and each integer $j\ge0$, put
\[
 P_a=(\Sym V)^{\otimes a},\qquad
 R_j=\bigl((\Sym^j V)^{\otimes a}\bigr)^{\Sa},\qquad
 R=\bigoplus_{j\ge0}R_j.
\]
The group $\Sa$ permutes the tensor factors, and multiplication in $P_a$
makes $R$ a graded commutative algebra. We use the averaging isomorphism
from a symmetric power to the corresponding invariant tensors: for any
complex vector space $W$, it sends $w_1\cdots w_a\in\Sym^a W$ to
\begin{equation}\label{eq:average}
 \Av(w_1\otimes\cdots\otimes w_a)
 =\frac1{a!}\sum_{\pi\in\Sa}
 w_{\pi(1)}\otimes\cdots\otimes w_{\pi(a)}.
\end{equation}
This map is well defined on the symmetric-power quotient, and the
quotient projection restricted to invariant tensors is its inverse.
In particular, $R_1=\Sym^a V$ and $R_b=\Sym^a(\Sym^b V)$ under these
identifications. For every integer $b\ge0$, multiplication in $P_a$ induces
a map $\Sym^b R_1\longrightarrow R_b$. Under the averaging identifications,
this is the \emph{canonical Foulkes--Howe map} considered here:
\begin{equation}\label{eq:map}
 \mu_{a,b,V}:\Sym^b(\Sym^a V)\longrightarrow\Sym^a(\Sym^b V).
\end{equation}
If $A\subseteq R$ is the subalgebra generated by $R_1$, then
$\operatorname{im}\mu_{a,b,V}=A_b$. In positive dimension, Raicu--Sam--Weyman
\cite[Section~2]{RSW} formulate this map as
$\Sym^b(D^aV)\longrightarrow D^a(\Sym^bV)$, with divided powers realized
as invariant tensors. Transporting both sides by the averaging above gives
exactly \eqref{eq:map}.

For fixed $a$ and $V$, the stabilization question asks for a nonnegative degree $b_0$
such that $\mu_{a,b,V}$ is surjective for every $b\ge b_0$. One may ask for
the least such degree in a fixed dimension, or for a bound valid in every
finite dimension at once. We prove the following uniform sufficient bound.

\begin{theorem}\label{thm:main}
For every integer $a\ge2$, every integer $b\ge a(a-1)$, and every
finite-dimensional complex vector space $V$, the map
$\mu_{a,b,V}$ in \eqref{eq:map} is surjective. Equivalently, $A_b=R_b$.
\end{theorem}

The case $a=1$ is the identity map in every degree.

\paragraph{History and significance.}
The affine Chow variety of decomposable $a$-forms consists of products of
$a$ linear forms in $\Sym^a(V^*)$. Its coordinate ring identifies with
$A$, whose normalization is $R$; here $b$ is the coordinate degree
\cite[Section~2, Theorem~2.3]{RSW}. This interpretation connects the canonical map to
the geometry of symmetric products; see also
Landsberg \cite[Section~7.6, Propositions~7.11--7.12]{Landsberg}.
Brion established eventual surjectivity \cite{Brion1993} and subsequently
gave an effective bound involving both $a$ and the dimension
\cite[Theorem~3.3]{Brion}. In the normalization picture, eventual
surjectivity says that the graded quotient $R/A$ has only finitely many
nonzero components. The least stabilization degree is one more than its
highest nonzero degree, or zero if $R=A$. Raicu, Sam, and Weyman bound the
regularity of the Segre ring and modules of covariants, including the
normalization \cite[Theorem~3.9]{RSW}. Their theorem does not compute the
top degree of the separate quotient $R/A$ that controls stabilization.
Theorem~\ref{thm:main} gives the sufficient bound $a(a-1)$ in every dimension.
It also answers positively the request for a polynomial bound in
\cite[Problem~7.19]{Landsberg}: in that source's notation $d=b$, $n=a$,
and $w=\dim V$, the bound is $d\ge n(n-1)$, including when $w=n$.
The present argument does not determine the least stabilization degree
in general; the bound is not asserted to be sharp.

The canonical-map question must be distinguished from the ordinary
Foulkes conjecture \cite{Foulkes1950}, which asks, for $a\le b$, for a
$\mathrm{GL}(V)$-equivariant injection
$\Sym^a(\Sym^b V)\hookrightarrow\Sym^b(\Sym^a V)$
\cite[Conjecture~2.10]{RSW}.
Surjectivity in Theorem~\ref{thm:main} implies that representation
comparison in the stated range: complete reducibility over $\C$ splits
the surjection and supplies an equivariant embedding in the opposite
direction. Corollary~\ref{cor:comparison} gives the precise statement.
Existence of an embedding alone does not specify the canonical map.
At $a=b$ the two representations are the same, so their ordinary comparison
is automatic, while a specified canonical endomorphism can still fail to be
invertible. Indeed, the canonical map has nonzero kernels at the diagonal pairs
$(a,b)=(5,5)$ and $(6,6)$ in suitable dimensions
\cite[arXiv version, Theorem~6(b),(c) and Remark~7]{CIM}, both outside the range of
Theorem~\ref{thm:main}.

\paragraph{The proof mechanism.}
An annihilator criterion reduces Theorem~\ref{thm:main} to a vanishing
statement for symmetric $a$-linear forms on $\Sym^b V$ whose diagonal
vanishes on products of $b$ vectors. Factoring a binary form first gives an
identity with a common product in every slot. Each of two differential
sequences then starts from an evaluation whose vanishing is sought and
inserts factors until it reaches an identity already known to be zero.
Injectivity recovers vanishing before those insertions. The first sequence
frees the last slot from the common product; induction lets the product
factors in the other slots vary independently; and the second sequence
removes their remaining common power from the vanishing conclusion.
A directional derivative in a vector variable $U$ toward a second variable
$Y$ lowers the $U$-degree and raises the $Y$-degree by one. The differential
lemma proves injectivity on a homogeneous space whenever the $U$-degree
exceeds the $Y$-degree. The first sequence uses a different $Y$ at each step;
the second raises the degree of one shared $Y$ at every step, and this
accumulation leads to the quadratic sufficient bound.

Successive weight shifts also occur in Ikenmeyer's proof
\cite[Section~5]{Ikenmeyer} of McKay's propagation theorem \cite{McKay}.
In Ikenmeyer's formulation, the canonical maps in the opposite direction
$\Sym^a(\Sym^b V)\longrightarrow\Sym^b(\Sym^a V)$ are considered with
$\dim V\ge b$ for each $b$. For fixed $a$ and $B\ge a$, injectivity at
$(a,B)$ propagates to every $(a,b)$ with $b\ge B$
\cite[Theorem~1.3 and Section~3]{Ikenmeyer}.
Here the two sequences surround an induction on the number of slots and
give a surjectivity range directly.
Section~\ref{sec:map} establishes the diagonal criterion,
Section~\ref{sec:differential} proves the one-step injectivity statement,
and Section~\ref{sec:vanishing} carries out the induction.

\section{The canonical image and diagonal vanishing}\label{sec:map}

We first express surjectivity of the specified multiplication map as a
vanishing problem. Fix integers $a,b\ge1$ and put $E=\Sym^b V$. An element
$P=u_1\cdots u_b\in E$ will be called a \emph{product of $b$ vectors};
the vectors may repeat or be zero.

For every integer $k\ge1$, pure powers span $\Sym^k V$. Indeed, the
polarization identity
\[
 k!\,u_1\cdots u_k
 =\sum_{J\subseteq\{1,\ldots,k\}}(-1)^{k-|J|}
   \left(\sum_{j\in J}u_j\right)^k
\]
holds in the symmetric algebra. On expanding first with the $u_i$ as formal
generators, every monomial missing an index cancels, and the only degree-$k$
monomial containing all $k$ indices has coefficient $k!$. Substituting
arbitrary vectors proves the identity. Products span $\Sym^k V$, and
$k!\ne0$ over $\C$, so pure powers span it as well.

Averaging sends $u^a\in\Sym^a V$ exactly to
$z(u)=u^{\otimes a}\in R_1$. The spanning fact with $k=a$ therefore says
that the elements $z(u)$ span $R_1$. Componentwise multiplication gives
\[
 z(u_1)\cdots z(u_b)=(u_1\cdots u_b)^{\otimes a},
\]
so multiplying spanning generators yields the exact image description
\[
 A_b=\operatorname{span}_{\C}
 \{P^{\otimes a}: P\text{ is a product of }b\text{ vectors}\}.
\]

For $\lambda\in R_b^*$, define the symmetric $a$-linear form on $E$
\[
 T_\lambda(w_1,\ldots,w_a)
 =\lambda\!\left(\Av(w_1\otimes\cdots\otimes w_a)\right).
\]
Every symmetric $a$-linear form arises this way: its linearization on
$E^{\otimes a}$ restricts to a functional on $R_b$, and symmetry makes
averaging leave its value unchanged. The functional is unique because
averages of simple tensors span $R_b$. For a product $P$,
$T_\lambda(P,\ldots,P)=\lambda(P^{\otimes a})$. Thus the annihilator of
$A_b$ corresponds exactly to the symmetric forms whose diagonal vanishes
on every product of $b$ vectors. Since $R_b$ is finite-dimensional,
$A_b=R_b$ if and only if that annihilator is zero. We have proved the
following criterion:

\medskip
\noindent
\emph{The map $\mu_{a,b,V}$ is surjective if and only if every symmetric
$a$-linear form $T$ on $\Sym^b V$ satisfying $T(P,\ldots,P)=0$ for all
products $P$ of $b$ vectors is zero.}
\medskip

This also explains the reversed roles of the two parameters in the equation
interpretation. The multiplication algebra above concerns products of $a$
linear forms in coordinate degree $b$. Here $T(P,\ldots,P)$ is a homogeneous
degree-$a$ polynomial on $\Sym^b V$, tested on products of $b$ vectors, viewed as linear
forms on $V^*$. The pairing just proved connects these two descriptions;
compare the product parametrization in \cite[Section~2]{RSW}.

\section{Differential injectivity}\label{sec:differential}

The induction will recover a polynomial from a sequence of its
directional derivatives. We first give the injectivity criterion that
permits each derivative to be cancelled.

Choose a basis of $V$, of dimension $n>0$. By a vector variable we mean an
independent block of $n$ scalar coordinate variables. For two such blocks
$U=(U_1,\ldots,U_n)$ and $Y=(Y_1,\ldots,Y_n)$, define
\[
 \D{Y}{U}=\sum_{k=1}^n Y_k\frac{\partial}{\partial U_k}.
\]
This operator lowers the $U$-degree by one and raises the $Y$-degree by one.
The criterion below is the positive-weight injectivity in the usual
$\mathfrak{sl}_2$ weight strings; compare
\cite[Section~11.1, Claim~11.4 and equation~(11.5)]{FultonHarris}.
Ikenmeyer's argument \cite[Section~5, Claim~5.1]{Ikenmeyer} decomposes
tensors according to untouched positions before applying a two-coordinate
weight shift.
We give a direct inner-product proof with all additional variables present.

\begin{lemma}\label{lem:shift}
Let $p,q$ be nonnegative integers. On polynomials homogeneous of degree $p$
in $U$ and degree $q$ in $Y$, the operator $\D{Y}{U}$ is injective if
$p>q$. The polynomials may depend on any finite collection of additional
variables independent of $U$ and $Y$.
\end{lemma}

\begin{proof}
Give the polynomial ring the positive definite Fischer, or apolar,
Hermitian inner product
\cite[Section~2, equation~(8) and Proposition~3]{AldazRender}, in which
distinct monomials are orthogonal and
\[
 \|X^\alpha\|^2=\alpha!:=\prod_k\alpha_k!;
\]
here $X$ lists all coordinates, including the additional variables.
Differentiation by a coordinate is adjoint to multiplication by that
coordinate. Thus, for $D=\D{Y}{U}$ and $E=\D{U}{Y}$, one has $D^*=E$.
A direct calculation gives
\[
 ED-DE=N_U-N_Y,
 \qquad
 N_U=\sum_k U_k\partial_{U_k},\quad
 N_Y=\sum_k Y_k\partial_{Y_k}.
\]
Consequently, for a polynomial $h$ of the indicated bidegree,
\begin{equation}\label{eq:norm}
 \|Dh\|^2-\|Eh\|^2=(p-q)\|h\|^2.
\end{equation}
If $Dh=0$ and $p>q$, the left side is nonpositive and the right side is
nonnegative, with equality only when $h=0$.
\end{proof}

Every application of Lemma~\ref{lem:shift} below concerns the entire
intermediate homogeneous space. Accordingly, a sequence of the operators
is injective whenever its degree inequality holds before each individual
operator.

\section{Vanishing on products}\label{sec:vanishing}

We now prove the vanishing statement in the criterion of
Section~\ref{sec:map}. Its formulation uses only symmetric multilinear
forms and products in $\Sym V$.

\begin{proposition}\label{prop:vanishing}
Let $V$ be a finite-dimensional complex vector space, and let $a\ge1$ and
$b\ge1$ be integers satisfying $b\ge a(a-1)$.
If a symmetric $a$-linear form $T$ on $\Sym^b V$ satisfies
\begin{equation}\label{eq:diagonal}
 T(P,\ldots,P)=0
 \qquad\text{for every product $P$ of $b$ vectors},
\end{equation}
then $T=0$.
\end{proposition}

\begin{proof}
The assertion is immediate for $V=0$, so suppose $\dim V>0$.
We induct on $a$, proving the assertion simultaneously for every admissible
$b$ and every finite-dimensional $V$. For $a=1$, products span $\Sym^b V$.
Now let $a\ge2$, assume the assertion for $a-1$, and put
\begin{equation}\label{eq:rm}
 r=a-1,\qquad m=b-r.
\end{equation}
Our hypotheses give
\begin{equation}\label{eq:range}
 b\ge r(r+1),\qquad m\ge r^2\ge1.
\end{equation}

\paragraph{Step 1: a split binary form.}
Let $Q=y_1\cdots y_m$ be a product of $m$ vectors. For every $c\in\C$
the homogeneous binary polynomial $X^r+cY^r$ splits into $r$ linear
factors over $\C$: for $c\ne0$, choose $\alpha\in\C$ with $\alpha^r=-c$ and
$\zeta=e^{2\pi i/r}$, so that
\[
 X^r+cY^r=\prod_{j=0}^{r-1}(X-\alpha\zeta^jY).
\]
For $c=0$ the factorization is $X^r$. Substitution in $\Sym V$ therefore
shows, without requiring $x$ and $t$ to be independent, that
\[
 P=(x^r+ct^r)Q
\]
is a product of $b$ vectors for all $x,t\in V$. For fixed $x,t,Q$,
condition \eqref{eq:diagonal} says that the polynomial $T(P,\ldots,P)$
in $c$ vanishes identically. By symmetry and multilinearity, its
coefficient of $c$ is
$a\,T(x^rQ,\ldots,x^rQ,t^rQ)$. Hence
\begin{equation}\label{eq:binary}
 T(\underbrace{x^rQ,\ldots,x^rQ}_{r\text{ slots}},t^rQ)=0.
\end{equation}

\paragraph{Step 2: freeing the last slot.}
We seek the same vanishing with the last input replaced by the independent
pure power $t^b$. Differentiation will insert the missing factors of $Q$
into that input, producing \eqref{eq:binary}; injectivity will then recover
the vanishing before those insertions. Fix $x\in V$, regard
$t,y_1,\ldots,y_m$ as independent vector variables, and set
\[
 H_x=T(\underbrace{x^rQ,\ldots,x^rQ}_{r\text{ slots}},t^b).
\]
It is homogeneous of degree $b$ in $t$ and degree $r$ in each $y_i$.
Differentiation in $t$ acts only on the last slot. Thus, for $0\le k\le m$,
\begin{equation}\label{eq:first-composition}
 \D{y_k}{t}\cdots\D{y_1}{t}H_x
 =\frac{b!}{(b-k)!}
 T(x^rQ,\ldots,x^rQ,t^{b-k}y_1\cdots y_k),
\end{equation}
with the empty composition understood when $k=0$.
At $k=m$ the right side is zero by \eqref{eq:binary}, first for every
choice of vector values and hence as a polynomial. Here and below we use
the elementary fact that a polynomial in finitely many variables over
$\C$ vanishing at every point is zero; it follows by induction on the
number of variables from the univariate root bound.
Immediately before the $i$-th operator, the degrees in $t$ and $y_i$ are
\[
 p=b-i+1,\qquad q=r,\qquad p-q=m-i+1\ge1.
\]
Here $y_i$ already has degree $r$ in $H_x$, but no earlier operator has
increased its degree. Lemma~\ref{lem:shift} makes every operator injective,
and hence $H_x=0$.

\paragraph{Step 3: induction on the remaining slots.}
For fixed $x,t\in V$, the expression
\[
 S_{x,t}(Q_1,\ldots,Q_r)
 :=T(x^rQ_1,\ldots,x^rQ_r,t^b)
\]
is a symmetric $r$-linear form on $\Sym^m V$. Step~2 says that its diagonal
vanishes at every product of $m$ vectors. Since
$m\ge r^2\ge r(r-1)$ and $m\ge1$, the induction hypothesis applies and gives
\begin{equation}\label{eq:inductive-vanishing}
 T(x^rQ_1,\ldots,x^rQ_r,t^b)=0
 \qquad(Q_1,\ldots,Q_r\in\Sym^m V).
\end{equation}
This identity holds for every $x,t\in V$.

\paragraph{Step 4: removing the common power.}
The factors $Q_i$ in \eqref{eq:inductive-vanishing} may now be chosen
independently; their remaining shared constraint is the multiplier $x^r$.
We will differentiate evaluations on separately chosen pure powers into this known
zero and again use injectivity to recover their vanishing. Fix $t\in V$ and
form the polynomial
\[
 f=T(v_1^b,\ldots,v_r^b,t^b)
\]
in independent vector variables $v_1,\ldots,v_r$. Introduce one further
independent vector variable $x$; initially $f$ has $x$-degree zero.
Apply $\D{x}{v_i}^{\,r}$ successively for $i=1,\ldots,r$. Each derivative
acts only on its own $v_i$ variable, so it does not differentiate the
previously inserted $x$ factors or any other slot. The combined result is
\begin{equation}\label{eq:second-composition}
 \D{x}{v_r}^{\,r}\cdots\D{x}{v_1}^{\,r}f
 =\left(\frac{b!}{m!}\right)^r
 T(x^rv_1^m,\ldots,x^rv_r^m,t^b)=0.
\end{equation}
The last equality follows from \eqref{eq:inductive-vanishing} for every
choice of the vector values, hence as a polynomial identity over $\C$.

All $r^2$ individual operators raise the degree of the same variable block
$x$. Before any one is applied, its degree is therefore at most $r^2-1$.
The source $v_i$ of that operator has previously been differentiated at
most $r-1$ times, so its degree is at least
$b-r+1=m+1\ge r^2+1$. Thus the source degree exceeds the destination degree
before every operator, with all other blocks as spectators.
Lemma~\ref{lem:shift} makes the composition injective, and
\eqref{eq:second-composition} implies $f=0$.
As $t$ is arbitrary, $T$ vanishes on every tuple of pure powers.

Pure powers span $\Sym^b V$ by the polarization argument in
Section~\ref{sec:map}. Multilinearity now yields $T=0$, completing the
induction.
\end{proof}

\begin{proof}[Proof of Theorem~\ref{thm:main}]
For $a\ge2$ and $b\ge a(a-1)$, Proposition~\ref{prop:vanishing} supplies
the vanishing assertion in the criterion of Section~\ref{sec:map}.
That criterion gives the surjectivity of $\mu_{a,b,V}$.
\end{proof}

\section{Representation comparison and scope}\label{sec:comparison}

\begin{corollary}\label{cor:comparison}
For every finite-dimensional complex vector space $V$, every integer
$a\ge2$, and every integer $b\ge a(a-1)$, there exists a
$\mathrm{GL}(V)$-equivariant embedding
\[
 \Sym^a(\Sym^b V)\hookrightarrow\Sym^b(\Sym^a V).
\]
In particular, this gives the sixth symmetric-power comparison for $b\ge30$.
\end{corollary}

\begin{proof}
Averaging and multiplication commute with the action of $\mathrm{GL}(V)$,
so the surjection $\mu_{a,b,V}$ of Theorem~\ref{thm:main} is equivariant.
Finite-dimensional polynomial representations over $\C$ are completely
reducible \cite[Chapter~I, Appendix~A, Sections~7--8]{Macdonald}.
Choose an invariant complement $M$ to its kernel. The restriction
$\mu_{a,b,V}|_M$ is an equivariant isomorphism onto
$\Sym^a(\Sym^b V)$. Its inverse followed by the inclusion of $M$ in
$\Sym^b(\Sym^a V)$ is the required embedding. This construction asserts
existence and does not identify the embedding with a prescribed canonical map.
\end{proof}

\begin{remark}
The differential lemma holds for every finite number of coordinates and
arbitrary additional variables, while the degree comparisons depend only
on $a$ and $b$. Together with the finite-dimensional annihilator argument,
this proves the same bound in every finite dimension, without a
stable-dimension reduction. The written proof uses positivity of the
Fischer Hermitian inner product, splitting of binary forms over $\C$, and
the nonzero characteristic-zero scalar factors in the coefficient and
differentiation identities.
\end{remark}

\begin{thebibliography}{99}
\bibitem{AldazRender}
J.~M.~Aldaz and H.~Render,
\href{https://doi.org/10.1007/s13324-023-00844-4}{\emph{A Fischer type decomposition theorem from the apolar inner product}},
Analysis and Mathematical Physics \textbf{13} (2023), article~91.

\bibitem{Brion1993}
M.~Brion,
\href{https://doi.org/10.1007/BF03026558}{\emph{Stable properties of plethysm: on two conjectures of Foulkes}},
Manuscripta Mathematica \textbf{80} (1993), no.~4, 347--371.

\bibitem{Brion}
M.~Brion,
\href{https://smf.emath.fr/publications/sur-certains-modules-gradues-associes-aux-produits-symetriques}{\emph{Sur certains modules gradu\'es associ\'es aux produits sym\'etriques}},
in \emph{Alg\`ebre non commutative, groupes quantiques et invariants}
(Reims, 1995), S\'eminaires et Congr\`es \textbf{2},
Soci\'et\'e Math\'ematique de France, Paris, 1997, 157--183.

\bibitem{CIM}
M.-W.~Cheung, C.~Ikenmeyer, and S.~Mkrtchyan,
\href{https://doi.org/10.1016/j.jsc.2016.09.002}{\emph{Symmetrizing tableaux and the 5th case of the Foulkes conjecture}},
Journal of Symbolic Computation \textbf{80} (2017), part~3, 833--843.
\href{https://arxiv.org/abs/1509.03944v1}{arXiv:1509.03944v1}.

\bibitem{Foulkes1950}
H.~O.~Foulkes,
\href{https://doi.org/10.1112/jlms/s1-25.3.205}{\emph{Concomitants of the quintic and sextic up to degree four in the coefficients of the ground form}},
Journal of the London Mathematical Society \textbf{25} (1950), no.~3,
205--209.

\bibitem{FultonHarris}
W.~Fulton and J.~Harris,
\href{https://doi.org/10.1007/978-1-4612-0979-9}{\emph{Representation Theory: A First Course}},
Graduate Texts in Mathematics \textbf{129}, Springer-Verlag, New York, 1991.

\bibitem{Ikenmeyer}
C.~Ikenmeyer,
\emph{On McKay's propagation theorem for the Foulkes conjecture},
\href{https://arxiv.org/abs/1509.04957v1}{arXiv:1509.04957v1}, 2015.

\bibitem{Landsberg}
J.~M.~Landsberg,
\href{https://doi.org/10.1007/s11565-014-0202-7}{\emph{Geometric complexity theory: an introduction for geometers}},
Annali dell'Universit\`a di Ferrara \textbf{61} (2015), no.~1, 65--117.
\href{https://arxiv.org/abs/1305.7387}{arXiv:1305.7387}.
Locators refer to the published version.

\bibitem{Macdonald}
I.~G.~Macdonald,
\href{https://doi.org/10.1093/oso/9780198534891.001.0001}{\emph{Symmetric Functions and Hall Polynomials}},
second edition, Oxford Mathematical Monographs, Clarendon Press, Oxford, 1995.

\bibitem{McKay}
T.~McKay,
\href{https://doi.org/10.1016/j.jalgebra.2007.12.003}{\emph{On plethysm conjectures of Stanley and Foulkes}},
Journal of Algebra \textbf{319} (2008), no.~5, 2050--2071.

\bibitem{RSW}
C.~Raicu, S.~V.~Sam, and J.~Weyman,
\href{https://doi.org/10.1007/s10013-021-00527-2}{\emph{On some modules supported in the Chow variety}},
Vietnam Journal of Mathematics \textbf{50} (2022), no.~2, 501--521.
\href{https://arxiv.org/abs/2108.10910v1}{arXiv:2108.10910v1}.
Locators refer to this preprint version.
\end{thebibliography}

\end{document}
